\documentclass[11pt]{article}
\usepackage{amsmath,amsfonts}
\usepackage{graphicx}
%\usepackage{xcolor}
%\definecolor{gray}{rgb}{.75,.75,.75}
%\usepackage[landscape]{geometry}
%\usepackage{hyperref}
%\usepackage[bookmarks=true, pdfborder={0 0 .5}, urlbordercolor=gray]{hyperref}

\renewcommand{\thefootnote}{\fnsymbol{footnote}}


\addtolength{\topmargin}{-.9in}
\addtolength{\textheight}{1.6in}
\addtolength{\oddsidemargin}{-.4in}
\addtolength{\textwidth}{.8in}
\pagestyle{empty}

\newcommand{\ds}{\displaystyle}
\newcommand{\ts}{\textstyle}

\newcommand{\Z}{\mathbb{Z}}
\newcommand{\R}{\mathbb{R}}
\newcommand{\C}{\mathbb{C}}
\newcommand{\EE}{\mathcal{E}}
\newcommand{\tZ}{\tilde{\Z}}

\newcommand{\Sym}{\mathrm{Sym}}

\renewcommand{\circle}[1]{{\bigcirc\hspace{-.75em}#1}}

\begin{document}

\footnote[0]{Notes by Peter Magyar \ {\tt magyar@math.msu.edu}}
\vspace{-1em}

\centerline{\bf Math 133\hspace{1.5em} \hfill{\Large\bf
Integration by Parts}\hfill Stewart \S7.1}
\vspace{1em}

\noindent 
{\bf Review of integrals.}
The definite integral gives the cumulative total of many small parts, such as the slivers which add up to the area under a graph. Numerically, it is a limit of Riemann sums:
$$
\int_a^b\! f(x)\,dx \ =\
\lim_{n\to\infty}\,\sum_{i=1}^n f(x_i)\,\Delta x\,,
$$
where we divide the interval $x\in[a,b]$ into $n$ increments
of size $\Delta x=\frac{b-a}{n}$ with division points
$a<a{+}\Delta x<a{+}2\Delta x<\cdots < a{+}n\Delta x=b$, and $x_1,\ldots x_n$ are sample points from each increment.
This definition is not a theoretical curiosity: it is the reason integrals are relevant to physical problems, and it is the only way to evaluate most integrals: there is no algebraic way.

However, for sufficiently simple functions $f(x)$, we can evaluate integrals algebraically by the shortcut of the Second Fundamental Theorem of Calculus. 
This says that if $f(x)$ is the rate of change of some known antiderivative $F(x)$, then the integral of $f(x)$ is the cumulative total change of $F(x)$:
$$
F'(x)=f(x) 
\qquad\Longrightarrow\qquad 
\int_a^b f(x)\,dx \ =\ \left.F(x)\right|_{x=a}^{x=b}
\ =\ 
F(b)-F(a)\,.
$$
(The First Fundamental Theorem says that the definite integral gives an antiderivative even if there is no formula $F(x)$:
defining $I(x)= \int_a^x f(t)\,dt$, we have $I'(x)  = f(x)$.)

Algebraic integration is the process of finding antiderivative formulas, 
denoted as indefinite integrals $\int\! f(x)\,dx=F(x)+C$. 
The most direct method is to reverse Basic Derivatives, 
such as $(x^p)' = px^{p-1}$ reversing to $\int\! x^p\,dx =\frac{x^{p+1}}{p+1}$. 
Our only other integration method so far is the Substitution Method, 
which reverses the Chain Rule:
$$
\int\! f(g(x))\,g'(x)\,dx 
\ =\
\int\! f(u)\,du
\ =\ F(u) +C
\ =\ F(g(x)) + C,
$$
where $u=g(x)$, $du=g'(x)\,dx$, and $F'(u)=f(u)$.
\\[.5em]
{\bf Reversing the Product Rule.} Since we have:
$$
(f(x)\,g(x))' \ =\ 
f(x)\, g'(x) + g(x)\, f'(x)\,,
$$
we can take the antiderivative of both sides to give:
$$
f(x)\,g(x) \ =\ 
\int\! f(x)\,g'(x)\,dx + \int\! g(x)\,f'(x)\,dx\,,
$$
%or:
$$
\int\! f(x)\,g'(x)\,dx \ =\ f(x)\,g(x) - \int\! g(x)\,f'(x)\,dx\,.
$$
In Leibnitz notation, taking $u=f(x)$, $du=f'(x)\,dx$
and $v=g(x)$, $dv = g'(x)\,dx$:
$$
\int\! u\ dv \ =\ uv \,-\, \int\! v\ du\,.
$$
This method transforms the integral of a product $f(x)\,g'(x)$
into $f(x) \,g(x)$ minus the integral of $g(x)\,f'(x)$, the other term in the Product Rule; we can 
think of lowering $f(x)$ to its derivative $f'(x)$ and raising $g'(x)$ to its antiderivative $g(x)$.

\newpage

\noindent
{\bf Method for Integration by Parts.}

\begin{enumerate}

\item Given an indefinite integral $\int\! h(x)\,dx$, find a factor of the integrand $h(x)$ which you recognize as the derivative of a function $g(x)$: that is, write $h(x)=f(x)\cdot g'(x)$.

\item Taking $u=f(x)$, $dv = g'(x)\,dx$, transform the integral $\int\! u\,dv=\int\! f(x)\, g'(x)\,dx$ into $uv-\int\! v\,du \ =\ f(x)\,g(x)-\int\! g(x)\,f'(x)\,dx$. 

\item Simplify $g(x)\,f'(x)$, possibly using identities, and try to find its integral by other methods such as Substitution.

\item Sometimes you can repeat Steps 1 \& 2 on $\int\! g(x)\,f'(x)\,dx$  
with a \emph{different} $u,v$.\footnote{
Repeating with the {\it same} factorization $\int\! v\,du$ would
get back the original integral $\int\! u\,dv$.
}
This might further simplify integral so you can evaluate by other methods.

\item Instead of simplifying the integral, Step 3 or 4 might give an expression with the {\it same} integral you started with. Solve the resulting equation to find that integral.

\end{enumerate}
%
Notice that Step 1 is the same as for the Method of Substitution, where you want to find a factor of the integrand which is a known derivative $g'(x)$; but for Substitution, $g(x)$ should also appear as an inside function in the remaining factor: $h(x)=f(g(x))\cdot g'(x)$.
\\[1em]
{\sc example:} Evaluate $\int x\cos(x)\,dx$. There are two obvious candidates for $u,v$. 
First, if we take $u=\cos(x)$, $dv=x\,dx$, we get $du=-\sin(x)\,dx$, $v = \frac12 x^2$, and:
$$\begin{array}{ccccc}
\int u\,dv & = & uv & - & \int v\,du 
\\[.5em]
\int \cos(x)\,x\,dx & = & \cos(x)\left(\tfrac 12 x^2\right) & - & \int \tfrac12 x^2\,(-\sin(x))\,dx
\end{array}$$
Unfortunately, the new integral $\int\! x^2\sin(x)\,dx$ is harder than the original $\int\! x\cos(x)\,dx$. We must make a wiser choice of $u,v$, so that the derivative $du$ will be simpler than the original $u$, while the antiderivative $v$ will be no worse than the original $dv$.

The other obvious choice will work: take $u=x$, $dv = \cos(x)\,dx$, so that $du = 1\, dx$ and $v=\sin(x)$. Then:\footnote{
For brevity, we neglect the constant $+C$, 
though you should write it in formal work.
}
$$\begin{array}{ccccc}
\int u\,dv & = & uv & - & \int v\,du 
\\[.5em]
\int x\,\cos(x)\,dx & = & x\,\sin(x) & - & \int \sin(x)\,1\,dx
\\[.5em]
& = & x\sin(x) & + & \cos(x) .
\end{array}$$
Thus, Steps 1--3 were enough to integrate.

To check our answer, we reverse our Integration by Parts using the Product Rule:
$$
(x\sin(x) + \cos(x))' 
\ =\ 
x\sin'(x) + \sin(x)(x)' + \cos'(x)
$$
$$
\ =\
x\cos(x) + \sin(x) - \sin(x) \ =\ x\cos(x).
$$
\\[.5em]
{\sc example:} Evaluate  $\int x^2\,e^{-x}\,dx$. We should choose $u=x^2$, $dv=e^{-x}$, so that $du=2x\,dx$ is simpler, but $v = -e^{-x}$ is no more complicated:
$$\begin{array}{ccccc}
\int u\,dv & = & uv & - & \int v\,du 
\\[.5em]
\int x^2\,e^{-x}\,dx & = & x^2(-e^{-x}) & - & \int (-e^{-x})\,2x\,dx
\\[.5em]
&=& -x^2e^{-x}&+ &2\int\! x e^{-x}\,dx
\end{array}$$
Going on to Step 4, we repeat the process for the integral on the right side,
this time with $u = x$, $dv= e^{-x}\,dx$ and
$du = dx$, $v = -e^{-x}$:
$$\begin{array}{ccccc}
\int u\,dv & = & uv & - & \int v\,du 
\\[.5em]
\int x\,e^{-x}\,dx & = & x(-e^{-x}) & - & \int (-e^{-x})\,dx
\\[.5em]
&=& -xe^{-x}&+ &\int\! e^{-x}\,dx
\\[.5em]
&=& -xe^{-x}&+ &(-e^{-x})
\end{array}$$
\\[-1.5em]
Putting these together:
$$
\int x^2\,e^{-x}\,dx
\ =\
-x^2e^{-x}+ 2(-xe^{-x}-e^{-x})
\ =\ -(x^2+2x+2)e^{-x}.
$$
\\[-.2em]
\noindent
{\sc example:} Evaluate $\int e^x\sin(x)\,dx$.
Steps 1--4 give:
$$\begin{array}{cccccc}
\int e^x\sin(x)\,dx & = & e^x\sin(x) & - & \int e^x\cos(x)\,dx,
& u = \sin(x), v = e^x
\\[.5em]
& = & e^x\sin(x) & - & \left(e^x\cos(x) -
\int e^x(-\sin(x))\,dx\right),
& u = \cos(x), v = e^x.
\end{array}$$
We conclude:
$$
\int e^x\sin(x)\,dx \ =\ 
e^x\sin(x) -e^x\cos(x) -\int e^x\sin(x)\,dx.
$$
Since our integral $\int e^x\sin(x)\,dx$ appears on both sides, we go to Step 5 
and solve for it:
$$
\int e^x\sin(x)\,dx\ =\ 
\tfrac12\left(e^x\sin(x)-e^x\cos(x)\right).
$$
\\[-.2em]
\noindent 
{\sc example:} Evaluate $\int \ln(x)\,dx$.
Here there does not seem to be any $dv$ factor, but we can always take $dv=1\,dx$, so $v = x$:
$$\begin{array}{ccccc}
\int u\,dv & = & uv & - & \int v\,du 
\\[.5em]
\int \ln(x)\,1\,dx
&=&
\ln(x)\,x &-& \int x\,\frac1x\,dx
\\[.5em]
&=&
x\ln(x)&-&x.
\end{array}$$
\\[-0em]
\noindent 
{\sc example:} Evaluate $\int \sin^{-1}(x)\,dx$.\footnote{
Notation: $\sin^{-1}(x)=\arcsin(x)$, but $\sin(x)^{-1}=\frac{1}{\sin(x)}=\csc(x)$.
}
Again we must use $u=\sin^{-1}(x)$ and $dv = 1\,dx$, counting on the fact that $du$ is simpler than $u$:
$$\begin{array}{ccccc}
\int u\,dv & = & uv & - & \int v\,du 
\\[.5em]
\int \sin^{-1}(x)\,1\,dx
&=&
\sin^{-1}(x)\,x &-& \int x\,\frac1{\sqrt{1-x^2}}\,dx
\end{array}$$
Continuing Step 3, we use the substitution $z = 1-x^2$ on the right-hand integral:
$$\ts
\int x\,\frac1{\sqrt{1-x^2}}\,dx
\ =\
-\frac12 \int \frac1{\sqrt{1-x^2}}\,(-2x)\,dx
\ =\
-\frac12\int \frac1{\sqrt z}\,dz
\ =\
-\sqrt{z}
\ =\
-\sqrt{1{-}x^2}.
$$
Combining:
$$\ts
\int \sin^{-1}(x)\,dx
\ =\
x\sin^{-1}(x) - (-\sqrt{1{-}x^2})
\ =\
x\sin^{-1}(x) + \sqrt{1{-}x^2}.
$$



\end{document}

%%%%%%%%%%%  Picture  %%%%%%%%%%
\\[0em]
\centerline{
%\includegraphics[width=2in]
{1-8-Discont-iv.png}
\qquad \qquad 
%\includegraphics[width=2in]
{1-8-Discont-v.png} 
}
\vspace{0em}

\noindent
%%%%%%%%%%%%%%%%%%%%%%%%%

%%%%%%%%%%%  Table  %%%%%%%%%%
In fact, we get the following derivatives:
$$\begin{array}{|c||c|c|c|c|c|c|}
\hline &&&&&& \\[-.9em] 
f(x) & \sin(x) & \cos(x) & \tan(x)  & \sec(x) & \csc(x) & \cot(x)
\\[.2em] \hline &&&&&& \\[-.9em]
f'(x) &\cos(x) & -\sin(x) & \sec^2(x) & \tan(x)\sec(x) & -\cot(x)\csc(x) & -\csc^2(x)
\\[.1em] \hline
\end{array}$$
\\[1em]
%%%%%%%%%%%%%%%%%%%%%%%%%%



















